% =====================================================================
%  minggu-15.tex  --  Slide Pertemuan Minggu ke-15
%  Pengantar Kriptografi (MPFA014) -- Prodi Matematika FMIPA UBB
%
%  Topik silabus : Review CPMK-1 dan CPMK-2
%  Sumber materi : buku_ajar/chapters/01--14 (ringkasan empat belas
%                  pertemuan; Bab 15 dan Bab 16 tidak dipakai)
%  Sub-CPMK      : 1.1--1.4 dan 2.1--2.5
%
%  Kompilasi dari dalam folder ppt/:
%      compile-ppt.bat 15
% =====================================================================

\documentclass[aspectratio=169,11pt]{beamer}
% =====================================================================
%  beamer-preamble.tex  --  prasarana dan gaya slide kuliah
%
%  Mata kuliah : Pengantar Kriptografi (MPFA014)
%  Program     : Program Studi Matematika, FMIPA, Universitas Bale Bandung
%
%  Berkas ini TIDAK memuat \documentclass. Setiap berkas minggu memuatnya
%  sendiri lalu menarik berkas ini:
%
%      \documentclass[aspectratio=169,11pt]{beamer}
%      % =====================================================================
%  beamer-preamble.tex  --  prasarana dan gaya slide kuliah
%
%  Mata kuliah : Pengantar Kriptografi (MPFA014)
%  Program     : Program Studi Matematika, FMIPA, Universitas Bale Bandung
%
%  Berkas ini TIDAK memuat \documentclass. Setiap berkas minggu memuatnya
%  sendiri lalu menarik berkas ini:
%
%      \documentclass[aspectratio=169,11pt]{beamer}
%      % =====================================================================
%  beamer-preamble.tex  --  prasarana dan gaya slide kuliah
%
%  Mata kuliah : Pengantar Kriptografi (MPFA014)
%  Program     : Program Studi Matematika, FMIPA, Universitas Bale Bandung
%
%  Berkas ini TIDAK memuat \documentclass. Setiap berkas minggu memuatnya
%  sendiri lalu menarik berkas ini:
%
%      \documentclass[aspectratio=169,11pt]{beamer}
%      \input{beamer-preamble}
%      \title{...} \subtitle{...} \setWeekInfo{Minggu n}{Topik}
%
%  Berkas minggu hanya berisi \setWeekInfo dan lingkungan frame.
%  Kompilasi dijalankan DARI DALAM folder ppt/ (lihat compile-ppt.bat),
%  karena gambar buku dirujuk dengan jalur relatif ../buku_ajar/.
% =====================================================================

% ---------------------------------------------------------------- bahasa
\usepackage[utf8]{inputenc}
\usepackage[T1]{fontenc}
\usepackage[indonesian]{babel}
\usepackage{lmodern}
\usepackage{microtype}

% ------------------------------------------------------------- matematika
\usepackage{amsmath,amssymb,mathtools}

% ----------------------------------------------------------------- tabel
\usepackage{booktabs}
\usepackage{array}
\usepackage{tabularx}

% --------------------------------------------------------------- grafik
\usepackage{tikz}
\usetikzlibrary{arrows.meta,positioning,shapes.geometric,calc,fit,backgrounds}

% ---------------------------------------------------------- kode program
\usepackage{listings}
\lstdefinestyle{slidecode}{
    basicstyle=\ttfamily\scriptsize,
    keywordstyle=\color{ubbblue}\bfseries,
    commentstyle=\color{ubbgreen}\itshape,
    stringstyle=\color{ubbred},
    numbers=left,
    numberstyle=\tiny\color{gray},
    numbersep=6pt,
    frame=single,
    rulecolor=\color{ubbblue!40},
    backgroundcolor=\color{ubbblue!3},
    breaklines=true,
    showstringspaces=false,
    tabsize=4,
    columns=fullflexible,
    keepspaces=true,
}
\lstset{style=slidecode}

% =====================================================================
%  WARNA INSTITUSI
% =====================================================================
\definecolor{ubbblue}{RGB}{12, 62, 120}      % warna utama
\definecolor{ubbteal}{RGB}{0, 118, 110}      % aksen kedua
\definecolor{ubbgreen}{RGB}{22, 122, 74}     % untuk blok contoh
\definecolor{ubbred}{RGB}{168, 42, 42}       % untuk blok peringatan
\definecolor{ubbgray}{RGB}{70, 76, 84}       % teks sekunder

% =====================================================================
%  IDENTITAS MATA KULIAH
% =====================================================================
\newcommand{\CourseName}{Pengantar Kriptografi}
\newcommand{\CourseCode}{MPFA014}
\newcommand{\CourseSemester}{Semester Ganjil}
\newcommand{\CourseProgram}{Program Studi Matematika}
\newcommand{\CourseFaculty}{Fakultas Matematika dan Ilmu Pengetahuan Alam}
\newcommand{\CourseInstitution}{Universitas Bale Bandung}
\newcommand{\CourseLecturer}{Cecep Suwanda, S.Si., M.Kom.}

% Label minggu, disetel tiap berkas lewat \setWeekInfo
\newcommand{\WeekLabel}{}
\newcommand{\WeekTopic}{}
\newcommand{\setWeekInfo}[2]{%
    \renewcommand{\WeekLabel}{#1}%
    \renewcommand{\WeekTopic}{#2}%
}

% =====================================================================
%  TEMA: header, footer, warna blok
% =====================================================================
\usetheme{default}
\setbeamertemplate{navigation symbols}{}
\setbeamertemplate{blocks}[rounded][shadow=false]

\setbeamercolor{palette primary}{bg=ubbblue, fg=white}
\setbeamercolor{palette secondary}{bg=ubbgray, fg=white}
\setbeamercolor{palette tertiary}{bg=ubbteal, fg=white}

\setbeamercolor{structure}{fg=ubbblue}
\setbeamercolor{normal text}{fg=black!92}
\setbeamercolor{alerted text}{fg=ubbred}

\setbeamercolor{block title}{bg=ubbblue, fg=white}
\setbeamercolor{block body}{bg=ubbblue!6, fg=black!92}
\setbeamercolor{block title alerted}{bg=ubbred, fg=white}
\setbeamercolor{block body alerted}{bg=ubbred!6, fg=black!92}
\setbeamercolor{block title example}{bg=ubbgreen, fg=white}
\setbeamercolor{block body example}{bg=ubbgreen!6, fg=black!92}

\setbeamercolor{itemize item}{fg=ubbblue}
\setbeamercolor{itemize subitem}{fg=ubbteal}
\setbeamercolor{enumerate item}{fg=ubbblue}

\setbeamerfont{headline}{size=\scriptsize}
\setbeamerfont{footline}{size=\scriptsize}
\setbeamerfont{block title}{size=\normalsize}

% Badan blok rata kiri (bukan rata kanan-kiri): teks slide lebih mudah dibaca
% dan baris terakhir yang pendek tidak lagi memicu "Underfull \hbox".
\addtobeamertemplate{block body begin}{\raggedright}{}
\addtobeamertemplate{block alerted body begin}{\raggedright}{}
\addtobeamertemplate{block example body begin}{\raggedright}{}

% Header: nama mata kuliah dan kode MK
\setbeamertemplate{headline}{%
    \begin{beamercolorbox}[wd=\paperwidth,ht=2.6ex,dp=1.2ex,leftskip=1em,rightskip=1em]{palette primary}%
        \usebeamerfont{headline}\CourseName\ \textperiodcentered\ \CourseCode%
        \hfill \CourseProgram%
    \end{beamercolorbox}%
}

% Footer: minggu dan nomor halaman
\setbeamertemplate{footline}{%
    \begin{beamercolorbox}[wd=\paperwidth,ht=2.4ex,dp=1.2ex,leftskip=1em,rightskip=1em]{palette secondary}%
        \usebeamerfont{footline}\WeekLabel\ \textperiodcentered\ \WeekTopic%
        \hfill \insertframenumber\,/\,\inserttotalframenumber%
    \end{beamercolorbox}%
}

% Halaman judul: judul, subjudul, lalu blok identitas mata kuliah
\setbeamertemplate{title page}{%
    \begin{centering}
        \vskip1.0em
        {\usebeamerfont{title}\usebeamercolor[fg]{title}\inserttitle\par}
        \vskip0.3em
        {\usebeamerfont{subtitle}\usebeamercolor[fg]{subtitle}\insertsubtitle\par}
        \vskip1.2em
        \textcolor{ubbblue}{\rule{0.72\textwidth}{0.8pt}}
        \vskip1.2em
        {\small
            \textbf{\CourseName}\ (\CourseCode)\par
            \vskip0.4em
            \CourseProgram\par
            \CourseFaculty\par
            \CourseInstitution\par
            \vskip0.6em
            \CourseSemester\ \textperiodcentered\ \WeekLabel\par
            \vskip0.8em
            Dosen Pengampu: \CourseLecturer\par
        }
        \vskip1.0em
        \textcolor{ubbblue}{\rule{0.72\textwidth}{0.4pt}}
    \end{centering}
}

% =====================================================================
%  PERINTAH BANTU
% =====================================================================
\newcommand{\R}{\mathbb{R}}
\newcommand{\N}{\mathbb{N}}
\newcommand{\Z}{\mathbb{Z}}
\newcommand{\code}[1]{\texttt{#1}}

% Rujukan gambar dari buku: \bookfig[lebar]{NN}{namaberkas.png}
% Jalur relatif terhadap folder ppt/ (dikompilasi dari sana).
\newcommand{\bookfig}[3][0.80\linewidth]{%
    \includegraphics[width=#1]{../buku_ajar/chapters/#2/figures/#3}%
}

% Keterangan sumber gambar, sepadan dengan \sumber di buku
\newcommand{\figsumber}[1]{%
    \par\vspace{1pt}{\tiny\itshape Sumber: #1\par}%
}

% Penanda blok materi yang sering dipakai di slide
\newcommand{\slidejudul}[1]{{\large\bfseries\textcolor{ubbblue}{#1}}\par\vspace{2pt}}
\newcommand{\kunci}[1]{\textcolor{ubbred}{\textbf{#1}}}

%      \title{...} \subtitle{...} \setWeekInfo{Minggu n}{Topik}
%
%  Berkas minggu hanya berisi \setWeekInfo dan lingkungan frame.
%  Kompilasi dijalankan DARI DALAM folder ppt/ (lihat compile-ppt.bat),
%  karena gambar buku dirujuk dengan jalur relatif ../buku_ajar/.
% =====================================================================

% ---------------------------------------------------------------- bahasa
\usepackage[utf8]{inputenc}
\usepackage[T1]{fontenc}
\usepackage[indonesian]{babel}
\usepackage{lmodern}
\usepackage{microtype}

% ------------------------------------------------------------- matematika
\usepackage{amsmath,amssymb,mathtools}

% ----------------------------------------------------------------- tabel
\usepackage{booktabs}
\usepackage{array}
\usepackage{tabularx}

% --------------------------------------------------------------- grafik
\usepackage{tikz}
\usetikzlibrary{arrows.meta,positioning,shapes.geometric,calc,fit,backgrounds}

% ---------------------------------------------------------- kode program
\usepackage{listings}
\lstdefinestyle{slidecode}{
    basicstyle=\ttfamily\scriptsize,
    keywordstyle=\color{ubbblue}\bfseries,
    commentstyle=\color{ubbgreen}\itshape,
    stringstyle=\color{ubbred},
    numbers=left,
    numberstyle=\tiny\color{gray},
    numbersep=6pt,
    frame=single,
    rulecolor=\color{ubbblue!40},
    backgroundcolor=\color{ubbblue!3},
    breaklines=true,
    showstringspaces=false,
    tabsize=4,
    columns=fullflexible,
    keepspaces=true,
}
\lstset{style=slidecode}

% =====================================================================
%  WARNA INSTITUSI
% =====================================================================
\definecolor{ubbblue}{RGB}{12, 62, 120}      % warna utama
\definecolor{ubbteal}{RGB}{0, 118, 110}      % aksen kedua
\definecolor{ubbgreen}{RGB}{22, 122, 74}     % untuk blok contoh
\definecolor{ubbred}{RGB}{168, 42, 42}       % untuk blok peringatan
\definecolor{ubbgray}{RGB}{70, 76, 84}       % teks sekunder

% =====================================================================
%  IDENTITAS MATA KULIAH
% =====================================================================
\newcommand{\CourseName}{Pengantar Kriptografi}
\newcommand{\CourseCode}{MPFA014}
\newcommand{\CourseSemester}{Semester Ganjil}
\newcommand{\CourseProgram}{Program Studi Matematika}
\newcommand{\CourseFaculty}{Fakultas Matematika dan Ilmu Pengetahuan Alam}
\newcommand{\CourseInstitution}{Universitas Bale Bandung}
\newcommand{\CourseLecturer}{Cecep Suwanda, S.Si., M.Kom.}

% Label minggu, disetel tiap berkas lewat \setWeekInfo
\newcommand{\WeekLabel}{}
\newcommand{\WeekTopic}{}
\newcommand{\setWeekInfo}[2]{%
    \renewcommand{\WeekLabel}{#1}%
    \renewcommand{\WeekTopic}{#2}%
}

% =====================================================================
%  TEMA: header, footer, warna blok
% =====================================================================
\usetheme{default}
\setbeamertemplate{navigation symbols}{}
\setbeamertemplate{blocks}[rounded][shadow=false]

\setbeamercolor{palette primary}{bg=ubbblue, fg=white}
\setbeamercolor{palette secondary}{bg=ubbgray, fg=white}
\setbeamercolor{palette tertiary}{bg=ubbteal, fg=white}

\setbeamercolor{structure}{fg=ubbblue}
\setbeamercolor{normal text}{fg=black!92}
\setbeamercolor{alerted text}{fg=ubbred}

\setbeamercolor{block title}{bg=ubbblue, fg=white}
\setbeamercolor{block body}{bg=ubbblue!6, fg=black!92}
\setbeamercolor{block title alerted}{bg=ubbred, fg=white}
\setbeamercolor{block body alerted}{bg=ubbred!6, fg=black!92}
\setbeamercolor{block title example}{bg=ubbgreen, fg=white}
\setbeamercolor{block body example}{bg=ubbgreen!6, fg=black!92}

\setbeamercolor{itemize item}{fg=ubbblue}
\setbeamercolor{itemize subitem}{fg=ubbteal}
\setbeamercolor{enumerate item}{fg=ubbblue}

\setbeamerfont{headline}{size=\scriptsize}
\setbeamerfont{footline}{size=\scriptsize}
\setbeamerfont{block title}{size=\normalsize}

% Badan blok rata kiri (bukan rata kanan-kiri): teks slide lebih mudah dibaca
% dan baris terakhir yang pendek tidak lagi memicu "Underfull \hbox".
\addtobeamertemplate{block body begin}{\raggedright}{}
\addtobeamertemplate{block alerted body begin}{\raggedright}{}
\addtobeamertemplate{block example body begin}{\raggedright}{}

% Header: nama mata kuliah dan kode MK
\setbeamertemplate{headline}{%
    \begin{beamercolorbox}[wd=\paperwidth,ht=2.6ex,dp=1.2ex,leftskip=1em,rightskip=1em]{palette primary}%
        \usebeamerfont{headline}\CourseName\ \textperiodcentered\ \CourseCode%
        \hfill \CourseProgram%
    \end{beamercolorbox}%
}

% Footer: minggu dan nomor halaman
\setbeamertemplate{footline}{%
    \begin{beamercolorbox}[wd=\paperwidth,ht=2.4ex,dp=1.2ex,leftskip=1em,rightskip=1em]{palette secondary}%
        \usebeamerfont{footline}\WeekLabel\ \textperiodcentered\ \WeekTopic%
        \hfill \insertframenumber\,/\,\inserttotalframenumber%
    \end{beamercolorbox}%
}

% Halaman judul: judul, subjudul, lalu blok identitas mata kuliah
\setbeamertemplate{title page}{%
    \begin{centering}
        \vskip1.0em
        {\usebeamerfont{title}\usebeamercolor[fg]{title}\inserttitle\par}
        \vskip0.3em
        {\usebeamerfont{subtitle}\usebeamercolor[fg]{subtitle}\insertsubtitle\par}
        \vskip1.2em
        \textcolor{ubbblue}{\rule{0.72\textwidth}{0.8pt}}
        \vskip1.2em
        {\small
            \textbf{\CourseName}\ (\CourseCode)\par
            \vskip0.4em
            \CourseProgram\par
            \CourseFaculty\par
            \CourseInstitution\par
            \vskip0.6em
            \CourseSemester\ \textperiodcentered\ \WeekLabel\par
            \vskip0.8em
            Dosen Pengampu: \CourseLecturer\par
        }
        \vskip1.0em
        \textcolor{ubbblue}{\rule{0.72\textwidth}{0.4pt}}
    \end{centering}
}

% =====================================================================
%  PERINTAH BANTU
% =====================================================================
\newcommand{\R}{\mathbb{R}}
\newcommand{\N}{\mathbb{N}}
\newcommand{\Z}{\mathbb{Z}}
\newcommand{\code}[1]{\texttt{#1}}

% Rujukan gambar dari buku: \bookfig[lebar]{NN}{namaberkas.png}
% Jalur relatif terhadap folder ppt/ (dikompilasi dari sana).
\newcommand{\bookfig}[3][0.80\linewidth]{%
    \includegraphics[width=#1]{../buku_ajar/chapters/#2/figures/#3}%
}

% Keterangan sumber gambar, sepadan dengan \sumber di buku
\newcommand{\figsumber}[1]{%
    \par\vspace{1pt}{\tiny\itshape Sumber: #1\par}%
}

% Penanda blok materi yang sering dipakai di slide
\newcommand{\slidejudul}[1]{{\large\bfseries\textcolor{ubbblue}{#1}}\par\vspace{2pt}}
\newcommand{\kunci}[1]{\textcolor{ubbred}{\textbf{#1}}}

%      \title{...} \subtitle{...} \setWeekInfo{Minggu n}{Topik}
%
%  Berkas minggu hanya berisi \setWeekInfo dan lingkungan frame.
%  Kompilasi dijalankan DARI DALAM folder ppt/ (lihat compile-ppt.bat),
%  karena gambar buku dirujuk dengan jalur relatif ../buku_ajar/.
% =====================================================================

% ---------------------------------------------------------------- bahasa
\usepackage[utf8]{inputenc}
\usepackage[T1]{fontenc}
\usepackage[indonesian]{babel}
\usepackage{lmodern}
\usepackage{microtype}

% ------------------------------------------------------------- matematika
\usepackage{amsmath,amssymb,mathtools}

% ----------------------------------------------------------------- tabel
\usepackage{booktabs}
\usepackage{array}
\usepackage{tabularx}

% --------------------------------------------------------------- grafik
\usepackage{tikz}
\usetikzlibrary{arrows.meta,positioning,shapes.geometric,calc,fit,backgrounds}

% ---------------------------------------------------------- kode program
\usepackage{listings}
\lstdefinestyle{slidecode}{
    basicstyle=\ttfamily\scriptsize,
    keywordstyle=\color{ubbblue}\bfseries,
    commentstyle=\color{ubbgreen}\itshape,
    stringstyle=\color{ubbred},
    numbers=left,
    numberstyle=\tiny\color{gray},
    numbersep=6pt,
    frame=single,
    rulecolor=\color{ubbblue!40},
    backgroundcolor=\color{ubbblue!3},
    breaklines=true,
    showstringspaces=false,
    tabsize=4,
    columns=fullflexible,
    keepspaces=true,
}
\lstset{style=slidecode}

% =====================================================================
%  WARNA INSTITUSI
% =====================================================================
\definecolor{ubbblue}{RGB}{12, 62, 120}      % warna utama
\definecolor{ubbteal}{RGB}{0, 118, 110}      % aksen kedua
\definecolor{ubbgreen}{RGB}{22, 122, 74}     % untuk blok contoh
\definecolor{ubbred}{RGB}{168, 42, 42}       % untuk blok peringatan
\definecolor{ubbgray}{RGB}{70, 76, 84}       % teks sekunder

% =====================================================================
%  IDENTITAS MATA KULIAH
% =====================================================================
\newcommand{\CourseName}{Pengantar Kriptografi}
\newcommand{\CourseCode}{MPFA014}
\newcommand{\CourseSemester}{Semester Ganjil}
\newcommand{\CourseProgram}{Program Studi Matematika}
\newcommand{\CourseFaculty}{Fakultas Matematika dan Ilmu Pengetahuan Alam}
\newcommand{\CourseInstitution}{Universitas Bale Bandung}
\newcommand{\CourseLecturer}{Cecep Suwanda, S.Si., M.Kom.}

% Label minggu, disetel tiap berkas lewat \setWeekInfo
\newcommand{\WeekLabel}{}
\newcommand{\WeekTopic}{}
\newcommand{\setWeekInfo}[2]{%
    \renewcommand{\WeekLabel}{#1}%
    \renewcommand{\WeekTopic}{#2}%
}

% =====================================================================
%  TEMA: header, footer, warna blok
% =====================================================================
\usetheme{default}
\setbeamertemplate{navigation symbols}{}
\setbeamertemplate{blocks}[rounded][shadow=false]

\setbeamercolor{palette primary}{bg=ubbblue, fg=white}
\setbeamercolor{palette secondary}{bg=ubbgray, fg=white}
\setbeamercolor{palette tertiary}{bg=ubbteal, fg=white}

\setbeamercolor{structure}{fg=ubbblue}
\setbeamercolor{normal text}{fg=black!92}
\setbeamercolor{alerted text}{fg=ubbred}

\setbeamercolor{block title}{bg=ubbblue, fg=white}
\setbeamercolor{block body}{bg=ubbblue!6, fg=black!92}
\setbeamercolor{block title alerted}{bg=ubbred, fg=white}
\setbeamercolor{block body alerted}{bg=ubbred!6, fg=black!92}
\setbeamercolor{block title example}{bg=ubbgreen, fg=white}
\setbeamercolor{block body example}{bg=ubbgreen!6, fg=black!92}

\setbeamercolor{itemize item}{fg=ubbblue}
\setbeamercolor{itemize subitem}{fg=ubbteal}
\setbeamercolor{enumerate item}{fg=ubbblue}

\setbeamerfont{headline}{size=\scriptsize}
\setbeamerfont{footline}{size=\scriptsize}
\setbeamerfont{block title}{size=\normalsize}

% Badan blok rata kiri (bukan rata kanan-kiri): teks slide lebih mudah dibaca
% dan baris terakhir yang pendek tidak lagi memicu "Underfull \hbox".
\addtobeamertemplate{block body begin}{\raggedright}{}
\addtobeamertemplate{block alerted body begin}{\raggedright}{}
\addtobeamertemplate{block example body begin}{\raggedright}{}

% Header: nama mata kuliah dan kode MK
\setbeamertemplate{headline}{%
    \begin{beamercolorbox}[wd=\paperwidth,ht=2.6ex,dp=1.2ex,leftskip=1em,rightskip=1em]{palette primary}%
        \usebeamerfont{headline}\CourseName\ \textperiodcentered\ \CourseCode%
        \hfill \CourseProgram%
    \end{beamercolorbox}%
}

% Footer: minggu dan nomor halaman
\setbeamertemplate{footline}{%
    \begin{beamercolorbox}[wd=\paperwidth,ht=2.4ex,dp=1.2ex,leftskip=1em,rightskip=1em]{palette secondary}%
        \usebeamerfont{footline}\WeekLabel\ \textperiodcentered\ \WeekTopic%
        \hfill \insertframenumber\,/\,\inserttotalframenumber%
    \end{beamercolorbox}%
}

% Halaman judul: judul, subjudul, lalu blok identitas mata kuliah
\setbeamertemplate{title page}{%
    \begin{centering}
        \vskip1.0em
        {\usebeamerfont{title}\usebeamercolor[fg]{title}\inserttitle\par}
        \vskip0.3em
        {\usebeamerfont{subtitle}\usebeamercolor[fg]{subtitle}\insertsubtitle\par}
        \vskip1.2em
        \textcolor{ubbblue}{\rule{0.72\textwidth}{0.8pt}}
        \vskip1.2em
        {\small
            \textbf{\CourseName}\ (\CourseCode)\par
            \vskip0.4em
            \CourseProgram\par
            \CourseFaculty\par
            \CourseInstitution\par
            \vskip0.6em
            \CourseSemester\ \textperiodcentered\ \WeekLabel\par
            \vskip0.8em
            Dosen Pengampu: \CourseLecturer\par
        }
        \vskip1.0em
        \textcolor{ubbblue}{\rule{0.72\textwidth}{0.4pt}}
    \end{centering}
}

% =====================================================================
%  PERINTAH BANTU
% =====================================================================
\newcommand{\R}{\mathbb{R}}
\newcommand{\N}{\mathbb{N}}
\newcommand{\Z}{\mathbb{Z}}
\newcommand{\code}[1]{\texttt{#1}}

% Rujukan gambar dari buku: \bookfig[lebar]{NN}{namaberkas.png}
% Jalur relatif terhadap folder ppt/ (dikompilasi dari sana).
\newcommand{\bookfig}[3][0.80\linewidth]{%
    \includegraphics[width=#1]{../buku_ajar/chapters/#2/figures/#3}%
}

% Keterangan sumber gambar, sepadan dengan \sumber di buku
\newcommand{\figsumber}[1]{%
    \par\vspace{1pt}{\tiny\itshape Sumber: #1\par}%
}

% Penanda blok materi yang sering dipakai di slide
\newcommand{\slidejudul}[1]{{\large\bfseries\textcolor{ubbblue}{#1}}\par\vspace{2pt}}
\newcommand{\kunci}[1]{\textcolor{ubbred}{\textbf{#1}}}


\title{Satu Semester Kriptografi dalam Satu Peta}
\subtitle{Review CPMK-1 dan CPMK-2: dari layanan keamanan sampai steganografi}
\setWeekInfo{Minggu 15}{Review CPMK-1 dan CPMK-2}

\begin{document}

% ----------------------------------------------------------- halaman judul
\begin{frame}[plain]
    \titlepage
\end{frame}

% -------------------------------------------------------------------- peta
\begin{frame}{Peta Empat Belas Pertemuan}
    \scriptsize
    \vspace{-0.15cm}
    \begin{columns}[T,onlytextwidth]
        \begin{column}{0.49\textwidth}
            \begin{tabularx}{\linewidth}{@{}c >{\raggedright\arraybackslash}X@{}}
                \toprule
                \textbf{Mg} & \textbf{Topik dan bab} \\
                \midrule
                1 & Pengantar dan urgensi (Bab 1) \\
                \addlinespace[1pt]
                2 & Layanan keamanan dan serangan (Bab 2) \\
                \addlinespace[1pt]
                3 & Landasan matematika (Bab 3) \\
                \addlinespace[1pt]
                4 & Klasik I: substitusi dan transposisi (Bab 4 \S1, \S5) \\
                \addlinespace[1pt]
                5 & Klasik II: Vigen\`ere, Playfair, Hill, Enigma (Bab 4 \S2--\S4) \\
                \addlinespace[1pt]
                6 & One-time pad (Bab 4 \S6) \\
                \addlinespace[1pt]
                7 & Kriptografi modern dan cipher alir (Bab 5) \\
                \bottomrule
            \end{tabularx}
        \end{column}
        \begin{column}{0.49\textwidth}
            \begin{tabularx}{\linewidth}{@{}c >{\raggedright\arraybackslash}X@{}}
                \toprule
                \textbf{Mg} & \textbf{Topik dan bab} \\
                \midrule
                8 & \textbf{Ujian Tengah Semester} \\
                \addlinespace[1pt]
                9 & Cipher blok, mode, DES, Triple DES (Bab 6--7) \\
                \addlinespace[1pt]
                10 & AES dan cipher blok lain (Bab 8--9) \\
                \addlinespace[1pt]
                11 & Algoritma RSA (Bab 10) \\
                \addlinespace[1pt]
                12 & Diffie--Hellman dan ElGamal (Bab 11--12) \\
                \addlinespace[1pt]
                13 & Kriptografi kurva eliptik (Bab 14) \\
                \addlinespace[1pt]
                14 & Knapsack, steganografi, steganalisis (Bab 13; modul 19--20) \\
                \bottomrule
            \end{tabularx}
        \end{column}
    \end{columns}
    \vspace{0.15cm}
    \begin{block}{Referensi}
        \footnotesize
        Bab 1--14 buku ajar; Munir, IF4020 Kriptografi, ITB, 2026 (modul 01--18)
    \end{block}
    \vspace{-0.1cm}
    \begin{alertblock}{Dua pertemuan ujian}
        \footnotesize
        Pertemuan 15 adalah review; pertemuan 16 adalah UAS. Tidak ada materi
        baru pada pertemuan review ini.
    \end{alertblock}
\end{frame}

% ----------------------------------------------------- capaian pembelajaran
\begin{frame}{Capaian Pembelajaran Minggu 15}
    \scriptsize
    \vspace{-0.12cm}
    \begin{columns}[T,onlytextwidth]
        \begin{column}{0.49\textwidth}
            \begin{block}{CPMK-2: konsep dan urgensi}
                \begin{itemize}
                    \setlength{\itemsep}{0pt}
                    \item \textbf{2.1} Menjelaskan definisi, terminologi,
                          sejarah, dan empat layanan kriptografi (kerahasiaan,
                          integritas, autentikasi, dan non-repudiation).
                    \item \textbf{2.2} Menjelaskan urgensi kriptografi dalam
                          teknologi informasi berdasarkan kasus kebocoran data,
                          kesalahan umum penggunaan kriptografi, dan tantangan
                          masa depan.
                    \item \textbf{2.3} Menjelaskan landasan matematika dasar
                          untuk kriptografi, yaitu aritmetika modulo, relatif
                          prima, invers modulo, bilangan prima, serta entropi
                          dan teori informasi.
                    \item \textbf{2.4} Mengidentifikasi jenis-jenis serangan
                          terhadap kriptografi berdasarkan data yang tersedia,
                          sifat serangan pasif atau aktif, dan side-channel
                          attack.
                    \item \textbf{2.5} Menjelaskan gagasan steganografi dan
                          steganalisis sebagai pelengkap kriptografi dalam
                          penyembunyian informasi beserta urgensinya.
                \end{itemize}
            \end{block}
        \end{column}
        \begin{column}{0.49\textwidth}
            \begin{block}{CPMK-1: algoritma klasik dan modern}
                \begin{itemize}
                    \setlength{\itemsep}{0pt}
                    \item \textbf{1.1} Mengenal cipher klasik, yaitu substitusi
                          (Caesar, monoalfabetik, Affine), transposisi
                          (Rail Fence), Vigen\`ere, Playfair, Hill, dan Enigma,
                          beserta gagasan kriptanalisisnya (analisis frekuensi
                          dan metode Kasiski).
                    \item \textbf{1.2} Mengenal one-time pad serta stream cipher
                          modern, yaitu representasi bit, operasi XOR,
                          keystream generator, FSR, RC4, A5, Trivium, Salsa20,
                          dan ChaCha20.
                    \item \textbf{1.3} Mengenal block cipher dan mode
                          operasinya (ECB, CBC, CFB, OFB, CTR), jaringan
                          Feistel, DES, Triple DES, AES, serta beberapa block
                          cipher lain (GOST, RC5, RC6, Blowfish, Serpent,
                          Twofish, dan MARS).
                    \item \textbf{1.4} Mengenal algoritma kriptografi kunci
                          publik, yaitu RSA, Diffie--Hellman, ElGamal,
                          knapsack, dan Elliptic Curve Cryptography (ECC),
                          beserta persoalan matematis yang mendasarinya.
                \end{itemize}
            \end{block}
        \end{column}
    \end{columns}
\end{frame}

% =====================================================================
\section{Layanan, Serangan, dan Matematika}
% =====================================================================

\begin{frame}{Empat Layanan dan Bentuk Pelanggarannya}
    \scriptsize
    \begin{tabularx}{\linewidth}{@{}>{\raggedright\arraybackslash}p{2.0cm} >{\raggedright\arraybackslash}p{5.0cm} >{\raggedright\arraybackslash}X@{}}
        \toprule
        \textbf{Layanan} & \textbf{Pertanyaan yang dijawab} & \textbf{Bentuk pelanggaran} \\
        \midrule
        Kerahasiaan & Apakah isi pesan tertutup bagi penyadap? & Intersepsi lalu lintas tanpa enkripsi; penyadapan kabel, elektromagnetik, dan akustik. \\
        \addlinespace[2pt]
        Integritas & Apakah pesan belum diubah di perjalanan? & Modifikasi nomor rekening atau nominal pada perintah pembayaran. \\
        \addlinespace[2pt]
        Autentikasi & Apakah pengirim benar-benar pihak yang diklaim? & Surel dengan alamat pengirim palsu; pemutaran ulang pesan sah. \\
        \addlinespace[2pt]
        Non-repudiation & Dapatkah pengirim menyangkal pesannya? & Sengketa transaksi tanpa bukti kriptografis yang mengikat pengirim. \\
        \bottomrule
    \end{tabularx}
    \vspace{0.2cm}
    \begin{alertblock}{Prinsip Kerckhoffs}
        \footnotesize
        Algoritma boleh dan harus publik; \kunci{hanya kunci yang rahasia}.
        Keamanan adalah pernyataan tentang kesulitan komputasi, bukan tentang
        ketidaktahuan penyerang.
    \end{alertblock}
\end{frame}

\begin{frame}{Mengapa Kriptografi Penting}
    \footnotesize
    \begin{columns}[T,onlytextwidth]
        \begin{column}{0.49\textwidth}
            \begin{block}{Tiga pemicu urgensi}
                \begin{itemize}
                    \setlength{\itemsep}{2pt}
                    \item \textbf{Kebocoran data}: data pribadi dan keuangan
                          yang tersimpan tanpa enkripsi dapat dibaca utuh.
                    \item \textbf{Penyadapan komunikasi}: kabel, radiasi
                          elektromagnetik, dan getaran akustik membocorkan isi
                          percakapan.
                    \item \textbf{Transaksi elektronik}: tanda terima, perintah
                          bayar, dan identitas digital harus dapat dipercaya.
                \end{itemize}
            \end{block}
            \begin{exampleblock}{Kesalahan umum}
                \footnotesize
                Menganggap penyandian (\textit{encoding}) atau tokenisasi
                sebagai enkripsi; memakai kunci pendek; dan mengulang
                \textit{nonce} atau IV.
            \end{exampleblock}
        \end{column}
        \begin{column}{0.49\textwidth}
            \begin{block}{Tantangan masa depan}
                \begin{itemize}
                    \setlength{\itemsep}{2pt}
                    \item Kriptografi ringan untuk perangkat IoT.
                    \item Kriptografi pasca-kuantum.
                    \item Kriptografi berbasis kecerdasan buatan.
                    \item \textit{Selective encryption} untuk data multimedia.
                \end{itemize}
            \end{block}
            \begin{alertblock}{Bacaan untuk UAS}
                \footnotesize
                Urgensi bukan daftar berita, melainkan alasan mengapa layanan
                keamanan tidak dapat ditunda. Setiap kasus kebocoran menunjukkan
                satu layanan yang gagal.
            \end{alertblock}
        \end{column}
    \end{columns}
\end{frame}

\begin{frame}{Serangan: dari Informasi ke Tujuan}
    \scriptsize
    \begin{columns}[T,onlytextwidth]
        \begin{column}{0.49\textwidth}
            \begin{block}{Berdasarkan informasi yang dimiliki}
                \begin{enumerate}
                    \setlength{\itemsep}{1pt}
                    \item \textit{Ciphertext-only}: hanya cipherteks; yang
                          tersulit.
                    \item \textit{Known-plaintext}: pasangan plainteks dan
                          cipherteks.
                    \item \textit{Chosen-plaintext}: penyerang memilih
                          plainteks.
                    \item \textit{Adaptive chosen-plaintext}: pilihan berikutnya
                          bergantung pada jawaban sebelumnya.
                    \item \textit{Chosen-ciphertext}: penyerang memilih
                          cipherteks dan melihat hasil dekripsi.
                \end{enumerate}
            \end{block}
            \begin{exampleblock}{Dua contoh dari buku}
                \footnotesize
                Caesar \code{KHOORZRUOG} dipecahkan dengan 25 pergeseran.
                Vigen\`ere dengan plainteks \code{AAAAA} membocorkan kunci
                karena cipherteksnya sama dengan kunci.
            \end{exampleblock}
        \end{column}
        \begin{column}{0.49\textwidth}
            \begin{block}{Berdasarkan tujuan akhir}
                \begin{itemize}
                    \setlength{\itemsep}{1pt}
                    \item \textit{Key-recovery}: menemukan kunci.
                    \item \textit{Plaintext-recovery}: menemukan pesan.
                    \item \textit{Distinguishing}: membedakan cipherteks dari
                          acak.
                    \item \textit{Forgeries}: memalsukan pesan atau tag.
                    \item \textit{Replay}: mengirim ulang pesan sah.
                \end{itemize}
            \end{block}
            \begin{alertblock}{Keamanan komputasi}
                \footnotesize
                Algoritma disebut aman secara komputasi bila biaya serangan
                melampaui nilai informasi dan waktu serangan melampaui masa
                kerahasiaan informasi.
            \end{alertblock}
        \end{column}
    \end{columns}
\end{frame}

\begin{frame}{Pasif, Aktif, dan Saluran Samping}
    \scriptsize
    \begin{columns}[T,onlytextwidth]
        \begin{column}{0.52\textwidth}
            \begin{block}{Serangan pasif}
                Hanya menyadap tanpa mengubah data: intersepsi dan analisis
                lalu lintas. Serangan pasif \kunci{sulit dideteksi}.
            \end{block}
            \begin{block}{Serangan aktif}
                Mengintervensi komunikasi: interupsi, fabrikasi, modifikasi,
                pemutaran ulang, dan \textit{man-in-the-middle}.
            \end{block}
            \begin{alertblock}{Serangan saluran samping}
                Memanfaatkan kebocoran fisik: waktu eksekusi, konsumsi daya
                (SPA/DPA), radiasi elektromagnetik, dan penyuntikan kesalahan
                (DFA).
            \end{alertblock}
        \end{column}
        \begin{column}{0.46\textwidth}
            \centering
            \bookfig[0.80\linewidth]{02}{serangan-mitm.png}
            \figsumber{Bab 2, \textit{Pengantar Kriptografi}.}
            \vspace{0.1cm}
            \begin{exampleblock}{MITM}
                \footnotesize
                Penyerang meneruskan pesan sambil membacanya.
            \end{exampleblock}
        \end{column}
    \end{columns}
\end{frame}

\begin{frame}{Matematika yang Dipakai Berulang}
    \scriptsize
    \begin{tabularx}{\linewidth}{@{}>{\raggedright\arraybackslash}p{3.0cm} >{\raggedright\arraybackslash}p{4.5cm} >{\raggedright\arraybackslash}X@{}}
        \toprule
        \textbf{Konsep} & \textbf{Bentuk yang dipakai} & \textbf{Muncul kembali pada} \\
        \midrule
        Aritmetika modulo & $a = qn + r$, $0 \le r < n$; contoh $-24 \bmod 11 = 9$ & Hill ($\bmod 26$), AES ($\mathrm{GF}(2^8)$), semua aritmetika kunci publik \\
        \addlinespace[1pt]
        Kekongruenan & $a \equiv b \pmod n \iff n \mid (a-b)$ & RSA, ElGamal, Diffie--Hellman \\
        \addlinespace[1pt]
        PBB dan Euclid & $\operatorname{PBB}(a,b)=1$ berarti relatif prima & syarat invers pada Affine dan RSA \\
        \addlinespace[1pt]
        Invers modulo & $a \cdot a^{-1} \equiv 1 \pmod n$; contoh $4^{-1} \bmod 9 = 7$ & Affine, RSA, ElGamal, knapsack \\
        \addlinespace[1pt]
        Bilangan prima dan totient & $\varphi(pq)=(p-1)(q-1)$ & pembangkitan kunci RSA \\
        \addlinespace[1pt]
        Teorema Euler & $a^{\varphi(n)} \equiv 1 \pmod n$ bila $\operatorname{PBB}(a,n)=1$ & kebenaran dekripsi RSA \\
        \addlinespace[1pt]
        Akar primitif dan logaritma diskret & $b \equiv g^{x} \pmod p$; mencari $x$ sukar & Diffie--Hellman, ElGamal, ECC (ECDLP) \\
        \addlinespace[1pt]
        Entropi Shannon & $H(X)=-\sum_i p(x_i)\log_2 p(x_i)$ & analisis frekuensi dan penilaian kebocoran \\
        \bottomrule
    \end{tabularx}
\end{frame}

\begin{frame}{Entropi: Mengapa Bahasa Alamiah Rapuh}
    \begin{columns}[T,onlytextwidth]
        \begin{column}{0.52\textwidth}
            \footnotesize
            Entropi maksimum alfabet 26 huruf adalah
            $\log_2 26 = 4{,}7004$ bit/karakter. Untuk 256 karakter ASCII,
            maksimumnya $\log_2 256 = 8$ bit/karakter.
            \vspace{0.12cm}
            \begin{exampleblock}{Contoh berita Kota Ternate}
                Entropi plainteks $H(P)=3{,}8988$ bit/karakter, yaitu
                $82{,}95\%$ dari maksimum. Setelah enkripsi,
                $H(C)=4{,}3035$ bit/karakter atau $91{,}55\%$ dari maksimum.
            \end{exampleblock}
            \begin{alertblock}{Redundansi}
                Selisih entropi teks alamiah dan maksimumnya adalah
                \kunci{redundansi}. Itulah bahan bakar analisis frekuensi.
                Entropi tinggi bukan bukti keamanan, melainkan alat diagnosis.
            \end{alertblock}
        \end{column}
        \begin{column}{0.46\textwidth}
            \centering
            \bookfig[1.0\linewidth]{03}{frekuensi-huruf-ternate.png}
            \figsumber{Bab 3, \textit{Pengantar Kriptografi}.}
        \end{column}
    \end{columns}
\end{frame}

% =====================================================================
\section{Kriptografi Klasik dan One-Time Pad}
% =====================================================================

\begin{frame}{Perbandingan Keluarga Cipher Klasik}
    \scriptsize
    \renewcommand{\arraystretch}{0.94}
    \begin{tabularx}{\linewidth}{@{}>{\raggedright\arraybackslash}p{1.8cm} >{\raggedright\arraybackslash}p{3.3cm} >{\raggedright\arraybackslash}p{3.4cm} >{\raggedright\arraybackslash}X@{}}
        \toprule
        \textbf{Keluarga} & \textbf{Satuan dan kunci} & \textbf{Kekuatan} & \textbf{Titik lemah} \\
        \midrule
        Caesar & per huruf; 26 pergeseran & rumus sederhana & dicoba 25 kali \\
        \addlinespace[0pt]
        Monoalfabetik & per huruf; $26!$ kunci & menyembunyikan identitas huruf & frekuensi huruf tetap \\
        \addlinespace[0pt]
        Affine & $C \equiv mP+b \pmod{26}$; 300 kunci & dua parameter & dua pasangan $P$--$C$ cukup \\
        \addlinespace[0pt]
        Vigen\`ere & kunci kata berulang & abjad-majemuk & metode Kasiski \\
        \addlinespace[0pt]
        Playfair & bigram; matriks $5\times5$ & mengolah pasangan huruf & 676 bigram terdistribusi \\
        \addlinespace[0pt]
        Hill & blok; matriks $K$ & difusi lokal & $\det(K)$ relatif prima 26; \textit{known-plaintext} \\
        \addlinespace[0pt]
        Enigma & rotor; posisi awal & periode sangat panjang & mesin \textit{Bombe} \\
        \addlinespace[0pt]
        Transposisi & posisi huruf; kunci permutasi & distribusi huruf utuh & bigram dan trigram terbaca \\
        \addlinespace[0pt]
        Super-enkripsi & substitusi lalu transposisi & dua lapis pertahanan & tetap cipher klasik \\
        \bottomrule
    \end{tabularx}
\end{frame}

\begin{frame}{Kriptanalisis Klasik: Tiga Perkakas}
    \begin{columns}[T,onlytextwidth]
        \begin{column}{0.50\textwidth}
            \scriptsize
            \begin{block}{Analisis frekuensi}
                Huruf yang paling sering pada cipherteks cenderung mewakili
                huruf yang paling sering pada bahasa aslinya. Bigram dan
                trigram memperkuat dugaan.
            \end{block}
            \begin{block}{Indeks kebetulan}
                Sekitar $0{,}065$ untuk cipher abjad-tunggal dan sekitar
                $0{,}038$ untuk teks acak atau cipher abjad-majemuk.
            \end{block}
            \begin{alertblock}{Kasiski dan Hill}
                Jarak kemunculan ulang kriptogram membocorkan panjang kunci
                Vigen\`ere. Hill dapat dipecahkan dengan pasangan
                plainteks--cipherteks yang diketahui.
            \end{alertblock}
        \end{column}
        \begin{column}{0.48\textwidth}
            \centering
            \bookfig[0.92\linewidth]{04}{tabel-frekuensi-huruf.png}
            \figsumber{Bab 4, \textit{Pengantar Kriptografi}.}
        \end{column}
    \end{columns}
\end{frame}

\begin{frame}{One-Time Pad: Janji Sempurna dan Harganya}
    \footnotesize
    \begin{columns}[T,onlytextwidth]
        \begin{column}{0.52\textwidth}
            OTP (Vernam 1917; dibuktikan Shannon 1949) adalah
            \kunci{satu-satunya cipher berkerahasiaan sempurna}:
            $H(P \mid C)=H(P)$.
            \vspace{0.12cm}
            \begin{block}{Tiga syarat}
                \begin{enumerate}
                    \setlength{\itemsep}{1pt}
                    \item kunci benar-benar acak;
                    \item panjang kunci sama dengan panjang pesan;
                    \item kunci hanya dipakai satu kali.
                \end{enumerate}
            \end{block}
            \begin{alertblock}{Kelemahan praktis}
                Pembangkitan dan distribusi kunci sepanjang pesan sukar.
                Pemakaian ulang kunci menghasilkan \textit{two-time pad}:
                $C_1 \oplus C_2 = P_1 \oplus P_2$, sehingga kunci saling
                menghapus.
            \end{alertblock}
        \end{column}
        \begin{column}{0.46\textwidth}
            \centering
            \bookfig[1.0\linewidth]{04}{contoh-one-time-pad.png}
            \figsumber{Bab 4, \textit{Pengantar Kriptografi}.}
        \end{column}
    \end{columns}
\end{frame}

\begin{frame}{Contoh Klasik yang Harus Bisa Diulang}
    \scriptsize
    \begin{columns}[T,onlytextwidth]
        \begin{column}{0.49\textwidth}
            \begin{exampleblock}{Caesar, $k=3$}
                \code{temui saya di jembatan merah nanti malam}
                \par\vspace{2pt}
                $\longrightarrow$ \code{WHPXL VDBD GL MHPEDWDQ PHUDK QDQWL PDODP}
                \par\vspace{2pt}
                Huruf \code{t} $=19$ menjadi $(19+3)\bmod 26=22=\code{W}$.
            \end{exampleblock}
            \begin{exampleblock}{Vigen\`ere, kunci \code{soreang}}
                \code{terimapesanangulaikambing}
                \par\vspace{2pt}
                $\longrightarrow$ \code{LSIMMNVWGRRAAMMZRMKNSTWEK}
                \par\vspace{2pt}
                Tiga huruf pertama: $(19+18)\bmod 26=11=\code{L}$,
                $(4+14)\bmod 26=18=\code{S}$,
                $(17+17)\bmod 26=8=\code{I}$.
            \end{exampleblock}
        \end{column}
        \begin{column}{0.49\textwidth}
            \begin{exampleblock}{One-time pad}
                Plainteks \code{onetimepad}, kunci \code{tbfrgfarhm}.
                \par\vspace{2pt}
                $\longrightarrow$ \code{HOJKOREGHP}
                \par\vspace{2pt}
                Huruf pertama: $(14+19)\bmod 26=7=\code{H}$.
            \end{exampleblock}
            \begin{alertblock}{Pola yang sama}
                Ketiga contoh memakai aritmetika modulo 26. Yang berubah hanya
                \emph{cara memilih pergeseran}: tetap, berulang, atau acak
                sepanjang pesan.
            \end{alertblock}
            \begin{block}{Catatan}
                Spasi dan tanda baca dibuang sebelum perhitungan, lalu
                dikembalikan setelah dekripsi.
            \end{block}
        \end{column}
    \end{columns}
\end{frame}

% =====================================================================
\section{Cipher Alir}
% =====================================================================

\begin{frame}{Perbandingan Cipher Alir}
    \scriptsize
    \vspace{-0.06cm}
    \renewcommand{\arraystretch}{0.94}
    \begin{tabularx}{\linewidth}{@{}>{\raggedright\arraybackslash}p{1.8cm} >{\raggedright\arraybackslash}p{3.7cm} >{\raggedright\arraybackslash}p{4.4cm} >{\raggedright\arraybackslash}X@{}}
        \toprule
        \textbf{Cipher} & \textbf{Kunci dan IV} & \textbf{Struktur} & \textbf{Catatan} \\
        \midrule
        LFSR $n$ bit & muatan awal $n$ bit & register geser, umpan balik XOR & periode maksimum $2^n-1$ \\
        \addlinespace[0pt]
        RC4 & kunci 40--2048 bit & larik $S$ 256 byte; KSA lalu PRGA & byte awal berkorelasi dengan kunci; tidak lagi dipakai \\
        \addlinespace[0pt]
        A5/1 & kunci sesi 64 bit & tiga LFSR 19, 22, 23 bit; kaidah mayoritas & frame GSM 228 bit setiap 4,6 ms; dipecahkan 1994 \\
        \addlinespace[0pt]
        Trivium & kunci 80 bit; IV 80 bit & tiga NLFSR 93, 84, 111 bit; gerbang AND & pemanasan 1152 detak \\
        \addlinespace[0pt]
        Salsa20 & kunci 256 bit; nonce 64 bit & ARX pada 16 kata 32 bit & satu blok \textit{keystream} 512 bit \\
        \addlinespace[0pt]
        ChaCha20 & kunci 256 bit & ARX; putaran kolom dan diagonal & difusi lebih baik daripada Salsa20 \\
        \bottomrule
    \end{tabularx}
    \vspace{0.08cm}
    \begin{block}{Pelajaran}
        \footnotesize
        Keamanan cipher alir ditentukan oleh ada tidaknya \kunci{hubungan
        statistik} antara kunci dan \textit{keystream}, bukan oleh panjang
        kunci saja.
    \end{block}
\end{frame}

\begin{frame}{Dari LFSR ke RC4}
    \begin{columns}[T,onlytextwidth]
        \begin{column}{0.50\textwidth}
            \scriptsize
            \begin{block}{LFSR}
                Setiap detak menggeser register satu posisi. Bit yang keluar
                menjadi \textit{keystream}, bit umpan balik mengisi posisi
                kosong.
            \end{block}
            \begin{exampleblock}{LFSR 4 bit}
                Muatan awal \code{1111}, umpan balik $f=b_4\oplus b_1$.
                Keluaran 15 detak: \code{111101011001000}. Setelah 15 detak
                register kembali ke \code{1111}.
            \end{exampleblock}
            \begin{block}{RC4}
                KSA mengacak larik $S$ 256 byte dengan kunci; PRGA menarik
                \textit{keystream} dari larik itu. Implementasi hati-hati
                membuang 256--512 byte awal.
            \end{block}
        \end{column}
        \begin{column}{0.48\textwidth}
            \centering
            \bookfig[1.0\linewidth]{05}{lfsr.png}
            \figsumber{Bab 5, \textit{Pengantar Kriptografi}.}
            \vspace{0.12cm}
            \begin{alertblock}{Batas pemakaian}
                \footnotesize
                RC4 tidak lagi aman karena byte awal \textit{keystream}
                berkorelasi dengan kunci.
            \end{alertblock}
        \end{column}
    \end{columns}
\end{frame}

\begin{frame}{A5/1, Trivium, dan Keluarga ARX}
    \scriptsize
    \begin{columns}[T,onlytextwidth]
        \begin{column}{0.49\textwidth}
            \begin{block}{A5/1}
                Tiga LFSR berukuran 19, 22, dan 23 bit. Kaidah mayoritas
                menentukan register mana yang bergeser. Satu frame GSM
                berukuran 228 bit setiap 4,6 milidetik. Serangan
                \textit{known-plaintext} Anderson (1994) memecahkannya.
            \end{block}
            \begin{block}{Trivium}
                Tiga NLFSR berukuran 93, 84, dan 111 bit; gerbang AND membuat
                strukturnya nirliniar. Pemanasan 1152 detak mencampur kunci
                dan IV.
            \end{block}
        \end{column}
        \begin{column}{0.49\textwidth}
            \begin{block}{Salsa20 dan ChaCha20}
                Keduanya memakai operasi ARX (\textit{add, rotate, xor}) pada
                matriks keadaan $4 \times 4$ berisi 16 kata 32 bit. Satu
                pemanggilan menghasilkan blok \textit{keystream} 512 bit.
            \end{block}
            \begin{exampleblock}{Mengapa dapat paralel}
                Setiap blok bergantung pada kunci, \textit{nonce}, dan pencacah.
                Karena tidak merantai blok sebelumnya, beberapa blok dapat
                dihitung serentak.
            \end{exampleblock}
            \begin{alertblock}{Perbedaan}
                ChaCha20 mengubah tata letak keadaan dan bentuk
                \textit{quarter round} agar difusinya lebih cepat.
            \end{alertblock}
        \end{column}
    \end{columns}
\end{frame}

% =====================================================================
\section{Cipher Blok sampai AES}
% =====================================================================

\begin{frame}{Prinsip Cipher Blok: Confusion dan Diffusion}
    \begin{columns}[T,onlytextwidth]
        \begin{column}{0.52\textwidth}
            \scriptsize
            \begin{block}{Dua ukuran}
                Ukuran \textbf{blok} menentukan luas difusi; ukuran
                \textbf{kunci} menentukan ketahanan terhadap pencarian
                menyeluruh. Blok 64 bit hanya boleh dipakai sekitar $2^{32}$
                blok per kunci.
            \end{block}
            \begin{block}{Confusion dan diffusion}
                Confusion menyembunyikan hubungan plainteks, cipherteks, dan
                kunci melalui substitusi nirlanjar, misalnya S-box DES
                $\code{110100}\to\code{1001}$. Diffusion menyebarkan pengaruh
                satu bit ke banyak bit melalui permutasi dan pergeseran.
            \end{block}
            \begin{alertblock}{Jaringan Feistel}
                $L_i=R_{i-1}$ dan $R_i=L_{i-1}\oplus f(R_{i-1},K_i)$.
                Struktur ini reversibel untuk sembarang fungsi $f$, sehingga
                dekripsi memakai urutan kunci terbalik.
            \end{alertblock}
        \end{column}
        \begin{column}{0.46\textwidth}
            \centering
            \bookfig[1.0\linewidth]{06}{jaringan-feistel.png}
            \figsumber{Bab 6, \textit{Pengantar Kriptografi}.}
            \vspace{0.1cm}
            \begin{exampleblock}{Cipher mainan}
                \footnotesize
                $(P \oplus K) \ll 1$ memperlihatkan enkripsi dan dekripsi yang
                saling membalik, tetapi sangat lemah karena hanya satu putaran.
            \end{exampleblock}
        \end{column}
    \end{columns}
\end{frame}

\begin{frame}{Mode Operasi: Pilihan yang Menentukan Keselamatan}
    \scriptsize
    \begin{tabularx}{\linewidth}{@{}>{\raggedright\arraybackslash}p{1.4cm} >{\raggedright\arraybackslash}p{5.1cm} >{\raggedright\arraybackslash}X@{}}
        \toprule
        \textbf{Mode} & \textbf{Rumus} & \textbf{Sifat dan pemakaian} \\
        \midrule
        ECB & $C_i = E_K(P_i)$ & Blok identik menghasilkan cipherteks identik; pola data bocor. \\
        \addlinespace[1pt]
        CBC & $C_i = E_K(P_i \oplus C_{i-1})$, $C_0=\mathrm{IV}$ & Pola berulang hilang; galat menjalar ke blok berikutnya. \\
        \addlinespace[1pt]
        CFB & $C_j = P_j \oplus \operatorname{MSB}_s(E_K(X_j))$ & Cipher blok dipakai seperti cipher alir; umpan balik dari cipherteks. \\
        \addlinespace[1pt]
        OFB & $C_j = P_j \oplus \operatorname{MSB}_s(E_K(X_j))$ & Umpan balik dari keluaran $E_K$; galat tidak menjalar. \\
        \addlinespace[1pt]
        CTR & $C_j = P_j \oplus E_K(T_j)$ & Paralel dan akses acak; \kunci{counter tidak boleh berulang}. \\
        \bottomrule
    \end{tabularx}
    \vspace{0.12cm}
    \begin{exampleblock}{Batas mode}
        \footnotesize
        Mode mengatur \emph{cara memakai} cipher blok, bukan kekuatan ciphernya.
        Tidak satu pun mode memperpanjang kunci efektif.
    \end{exampleblock}
\end{frame}

\begin{frame}{DES, Double DES, dan Triple DES}
    \begin{columns}[T,onlytextwidth]
        \begin{column}{0.53\textwidth}
            \scriptsize
            \begin{block}{Parameter DES}
                Blok 64 bit; kunci eksternal 64 bit, tetapi hanya
                \kunci{56 bit} yang dipakai karena 8 bit paritas. Jaringan
                Feistel 16 putaran. Ruang kunci $2^{56} =
                72.057.594.037.927.936$.
            \end{block}
            \begin{exampleblock}{Biaya pencarian menyeluruh}
                Pada laju $10^6$ kunci/detik: sekitar 2.284 tahun terburuk dan
                1.142 tahun rata-rata.
            \end{exampleblock}
            \begin{block}{Double dan Triple DES}
                Double DES tumbang oleh \textit{meet-in-the-middle} dengan
                sekitar $2^{56}$ penyimpanan. Triple DES EDE memberi 112 bit
                (dua kunci) atau 168 bit (tiga kunci), tetapi blok 64 bit tetap
                menjadi batasnya.
            \end{block}
        \end{column}
        \begin{column}{0.45\textwidth}
            \centering
            \bookfig[0.86\linewidth]{07}{skema-enkripsi-des.png}
            \figsumber{Bab 7, \textit{Pengantar Kriptografi}.}
            \vspace{0.1cm}
            \begin{alertblock}{Fungsi putaran}
                \footnotesize
                Ekspansi $32\to48$ bit, XOR dengan kunci putaran, delapan S-box
                $6\to4$ bit, lalu permutasi $P$.
            \end{alertblock}
        \end{column}
    \end{columns}
\end{frame}

\begin{frame}{AES: Standar yang Berlaku Sekarang}
    \begin{columns}[T,onlytextwidth]
        \begin{column}{0.50\textwidth}
            \scriptsize
            \begin{block}{Parameter}
                Blok tetap \kunci{128 bit}; kunci 128, 192, atau 256 bit;
                putaran 10, 12, atau 14. \texttt{State} berupa matriks
                $4 \times 4$ byte yang diisi kolom demi kolom.
            \end{block}
            \begin{block}{Alur satu putaran}
                Putaran awal hanya \textit{AddRoundKey}. Lalu $N_r-1$ putaran
                berisi \textit{SubBytes}, \textit{ShiftRows},
                \textit{MixColumns}, dan \textit{AddRoundKey}. Putaran terakhir
                tidak memakai \textit{MixColumns}.
            \end{block}
            \begin{exampleblock}{Contoh \textit{MixColumns}}
                Kolom $[\code{26}\ \code{7B}\ \code{BD}\ \code{43}]^{T}$
                menjadi $[\code{3F}\ \code{4F}\ \code{F9}\ \code{2A}]^{T}$.
            \end{exampleblock}
        \end{column}
        \begin{column}{0.48\textwidth}
            \centering
            \bookfig[0.96\linewidth]{08}{alur-enkripsi-aes.png}
            \figsumber{Bab 8, \textit{Pengantar Kriptografi}.}
            \vspace{0.1cm}
            \begin{alertblock}{Aritmetika dan keamanan}
                \footnotesize
                Operasi berada di $\mathrm{GF}(2^8)$ dengan polinom reduksi
                $x^8+x^4+x^3+x+1$. Ruang kunci AES-128 berisi
                $2^{128}=3{,}4\times10^{38}$ kemungkinan.
            \end{alertblock}
        \end{column}
    \end{columns}
\end{frame}

\begin{frame}{Perbandingan Cipher Blok}
    \scriptsize
    \renewcommand{\arraystretch}{0.94}
    \begin{tabularx}{\linewidth}{@{}>{\raggedright\arraybackslash}p{1.7cm} >{\raggedright\arraybackslash}p{3.0cm} c >{\raggedright\arraybackslash}X@{}}
        \toprule
        \textbf{Cipher} & \textbf{Blok / kunci} & \textbf{Putaran} & \textbf{Ciri dan status} \\
        \midrule
        DES & 64 / 56 bit & 16 & Feistel; dipecahkan 1998--1999 \\
        \addlinespace[0pt]
        Triple DES & 64 / 112--168 bit & 48 & EDE; kompatibel DES; blok 64 bit \\
        \addlinespace[0pt]
        AES & 128 / 128--256 bit & 10/12/14 & SPN; standar saat ini \\
        \addlinespace[0pt]
        GOST (Magma) & 64 / 256 bit & 32 & Feistel; delapan S-box $4\times4$ \\
        \addlinespace[0pt]
        RC5 & 32/64/128 / 0--2040 bit & 0--255 & Tiga parameter bebas $w$, $r$, $b$ \\
        \addlinespace[0pt]
        RC6 & 128 / 128--2040 bit & --- & Empat register; perkalian $x(2x+1)$ \\
        \addlinespace[0pt]
        Blowfish & 64 / 32--448 bit & 16 & Empat S-box dan larik $P$ \\
        \addlinespace[0pt]
        Serpent & 128 / 128--256 bit & 32 & SPN murni; delapan S-box $4\times4$ \\
        \addlinespace[0pt]
        Twofish & 128 / sampai 256 bit & 16 & Empat S-box $8\times8$ bergantung kunci \\
        \addlinespace[0pt]
        MARS & 128 / 128--448 bit & 32 & Feistel tak seimbang: $8+16+8$ \\
        \bottomrule
    \end{tabularx}
    \vspace{0.1cm}
    \begin{block}{Bacaan}
        \footnotesize
        Semua baris adalah cipher blok. Yang berubah adalah struktur putaran,
        ukuran blok, dan cara kunci putaran diturunkan; ukuran blok menentukan
        batas aman pemakaian satu kunci.
    \end{block}
\end{frame}

% =====================================================================
\section{Kriptografi Kunci-Publik}
% =====================================================================

\begin{frame}{Gagasan Kunci-Publik dan Kriptografi Hibrida}
    \footnotesize
    \begin{columns}[T,onlytextwidth]
        \begin{column}{0.50\textwidth}
            \begin{block}{Dua kunci}
                Kunci publik boleh disebarkan; kunci privat dijaga. Yang
                dienkripsi dengan kunci publik hanya dapat dibuka dengan kunci
                privat yang bersesuaian.
            \end{block}
            \begin{block}{Persoalan yang mendasari}
                \begin{itemize}
                    \setlength{\itemsep}{1pt}
                    \item RSA: memfaktorkan $n=pq$.
                    \item Diffie--Hellman dan ElGamal: logaritma diskret.
                    \item ECC: logaritma diskret pada kurva eliptik.
                    \item Knapsack: himpunan bagian dengan jumlah tertentu.
                \end{itemize}
            \end{block}
        \end{column}
        \begin{column}{0.48\textwidth}
            \begin{exampleblock}{Kriptografi hibrida}
                Kunci sesi $K$ dienkripsi dengan RSA,
                $C_K=E_{\mathrm{RSA},\mathrm{PK}_{\mathrm{Bob}}}(K)$.
                Pesan dienkripsi dengan AES,
                $C_M=E_{\mathrm{AES},K}(M)$. Penerima memulihkan $K$ lalu $M$.
            \end{exampleblock}
            \begin{alertblock}{Dua kelemahan yang harus diingat}
                \begin{itemize}
                    \setlength{\itemsep}{1pt}
                    \item RSA tanpa \textit{padding} bersifat deterministik.
                    \item Distribusi kunci publik rentan
                          \textit{man-in-the-middle}.
                \end{itemize}
            \end{alertblock}
        \end{column}
    \end{columns}
\end{frame}

\begin{frame}{RSA: Keamanan yang Bersandar pada Faktorisasi}
    \scriptsize
    \begin{columns}[T,onlytextwidth]
        \begin{column}{0.51\textwidth}
            \begin{block}{Pembangkitan kunci}
                Pilih prima $p,q$; hitung $n=pq$ dan
                $\varphi(n)=(p-1)(q-1)$; pilih $e$ dengan
                $\operatorname{PBB}(e,\varphi(n))=1$; hitung
                $d \equiv e^{-1} \pmod{\varphi(n)}$.
                \[
                    c=m^{e} \bmod n, \qquad m=c^{d} \bmod n.
                \]
            \end{block}
            \begin{exampleblock}{Contoh dari bab}
                $p=47$, $q=71$ memberi $n=3337$, $\varphi(n)=3220$,
                $e=79$, dan $d=1019$. Pesan $m=16$ menjadi
                $c=16^{79}\bmod 3337=1493$, lalu
                $1493^{1019}\bmod 3337=16$. Pesan panjang dipecah menjadi blok
                $0 \le m_i \le n-1$.
            \end{exampleblock}
        \end{column}
        \begin{column}{0.47\textwidth}
            \begin{alertblock}{Keamanan dan kelemahan}
                \begin{itemize}
                    \setlength{\itemsep}{1pt}
                    \item Keamanan terletak pada sulitnya memfaktorkan $n$.
                    \item Modulus 256 bit pecah dalam beberapa jam; 512 bit
                          dengan beberapa ratus komputer. Bahan kuliah
                          menyatakan $n$ di atas 1024 bit baru aman.
                    \item Tanpa OAEP, cipherteks deterministik dan dapat
                          diserang \textit{chosen-ciphertext}. Distribusi kunci
                          publik juga rentan \textit{man-in-the-middle}.
                    \item Algoritma Shor memfaktorkan $n$ dalam waktu
                          polinomial pada komputer kuantum.
                \end{itemize}
            \end{alertblock}
        \end{column}
    \end{columns}
\end{frame}

\begin{frame}{Diffie--Hellman: Menyepakati Rahasia di Saluran Terbuka}
    \scriptsize
    \begin{columns}[T,onlytextwidth]
        \begin{column}{0.50\textwidth}
            \begin{block}{Protokol}
                Parameter publik: prima $p$ dan akar primitif $g$.
                Alice mengirim $A=g^{a}\bmod p$; Bob mengirim
                $B=g^{b}\bmod p$. Keduanya memperoleh
                \[
                    K=B^{a}=A^{b}=g^{ab}\bmod p.
                \]
            \end{block}
            \begin{exampleblock}{Contoh dari bab}
                $p=353$, $g=3$, $a=97$, $b=233$ memberi
                $A=40$, $B=248$, dan $K=160$ dari kedua arah.
            \end{exampleblock}
            \begin{alertblock}{Serangan MITM}
                Mallory mengganti $g^{a}$ dan $g^{b}$ dengan $g^{m}$, sehingga
                ia memegang $K_1=g^{am}$ dan $K_2=g^{bm}$. Penangkalnya adalah
                otentikasi, bukan memperbesar $p$.
            \end{alertblock}
        \end{column}
        \begin{column}{0.48\textwidth}
            \centering
            \bookfig[0.92\linewidth]{11}{contoh-numerik-dh.png}
            \figsumber{Bab 11, \textit{Pengantar Kriptografi}.}
            \vspace{0.1cm}
            \begin{block}{Perluasan}
                \footnotesize
                Protokol dapat diperluas ke tiga pihak dengan pola melingkar;
                ketiganya sampai pada $K=g^{abc}\bmod p$.
            \end{block}
        \end{column}
    \end{columns}
\end{frame}

\begin{frame}{ElGamal: Enkripsi yang Memakai Kunci Sesaat}
    \scriptsize
    \begin{columns}[T,onlytextwidth]
        \begin{column}{0.50\textwidth}
            \begin{block}{Kunci dan enkripsi}
                Kunci privat $x$; kunci publik $y=g^{x}\bmod p$. Untuk pesan
                $m$ dan kunci sesaat acak $k$:
                \[
                    a=g^{k}\bmod p, \qquad b=y^{k}\cdot m \bmod p.
                \]
                Cipherteksnya pasangan $(a,b)$.
            \end{block}
            \begin{block}{Dekripsi}
                Hitung $(a^{x})^{-1} \equiv a^{\,p-1-x} \pmod p$, lalu
                $m=b\cdot(a^{x})^{-1}\bmod p$.
            \end{block}
            \begin{exampleblock}{Contoh dari bab}
                $p=2273$, $g=3$, $x=243$ memberi $y=461$. Pesan $m=700$
                dengan $k=1463$ menjadi $(1439,74)$ dan kembali menjadi $700$.
            \end{exampleblock}
        \end{column}
        \begin{column}{0.48\textwidth}
            \begin{block}{Yang menguntungkan}
                Karena $k$ dipilih acak untuk setiap pesan, plainteks yang sama
                dapat menghasilkan cipherteks berbeda.
            \end{block}
            \begin{alertblock}{Yang merugikan}
                \begin{itemize}
                    \setlength{\itemsep}{2pt}
                    \item Cipherteks dua kali panjang plainteks.
                    \item Setiap blok pesan harus berada di $[0,p-1]$.
                    \item Keamanan bergantung pada sulitnya logaritma diskret.
                \end{itemize}
            \end{alertblock}
            \begin{exampleblock}{Akar primitif}
                \footnotesize
                $a$ adalah akar primitif modulo $p$ bila
                $a^{1},a^{2},\dots,a^{p-1}$ menghasilkan seluruh residu bukan
                nol tepat satu kali.
            \end{exampleblock}
        \end{column}
    \end{columns}
\end{frame}

\begin{frame}{Perbandingan Keluarga Kunci-Publik}
    \scriptsize
    \begin{tabularx}{\linewidth}{@{}>{\raggedright\arraybackslash}p{1.6cm} >{\raggedright\arraybackslash}p{3.2cm} >{\raggedright\arraybackslash}p{3.5cm} >{\raggedright\arraybackslash}X@{}}
        \toprule
        \textbf{Algoritma} & \textbf{Kegunaan} & \textbf{Persoalan matematis} & \textbf{Kelemahan utama} \\
        \midrule
        RSA & Enkripsi dan pembungkusan kunci sesi & Memfaktorkan $n=pq$ & Lambat; deterministik tanpa OAEP; butuh modulus besar \\
        \addlinespace[1pt]
        Diffie--Hellman & Menyepakati kunci bersama & Logaritma diskret & Tidak saling mengotentikasi; rentan MITM \\
        \addlinespace[1pt]
        ElGamal & Enkripsi pesan atau kunci & Logaritma diskret & Cipherteks dua kali panjang; $k$ harus acak per pesan \\
        \bottomrule
    \end{tabularx}
    \vspace{0.12cm}
    \begin{block}{Benang merah}
        \footnotesize
        Ketiganya tidak menyembunyikan pesan dengan substitusi biasa, melainkan
        bersandar pada satu \kunci{persoalan komputasi yang diyakini sukar}.
        Karena itu panjang kunci dan struktur parameter menentukan keamanan.
    \end{block}
    \begin{alertblock}{Ancaman yang sama}
        \footnotesize
        Algoritma Shor memecahkan RSA, Diffie--Hellman, dan ElGamal pada komputer
        kuantum. Ketiganya harus bermigrasi ke skema pasca-kuantum.
    \end{alertblock}
\end{frame}

% =====================================================================
\section{ECC, Knapsack, dan Steganografi}
% =====================================================================

\begin{frame}{ECC: Keamanan pada Kunci yang Pendek}
    \scriptsize
    \begin{columns}[T,onlytextwidth]
        \begin{column}{0.50\textwidth}
            \begin{block}{Kurva dan grup}
                $y^{2}=x^{3}+ax+b$ dengan $4a^{3}+27b^{2}\neq0$. Titik-titiknya
                bersama titik di ketakhinggaan $\mathcal{O}$ membentuk grup
                abelian.
            \end{block}
            \begin{block}{Operasi titik}
                Penjumlahan memakai gradien
                $m=(y_p-y_q)(x_p-x_q)^{-1}$, penggandaan memakai
                $m=(3x_p^{2}+a)(2y_p)^{-1}$. Koordinat hasil:
                $x_r=m^{2}-x_p-x_q$ dan
                $y_r=m(x_p-x_r)-y_p$, semuanya modulo $p$.
            \end{block}
            \begin{alertblock}{ECDLP}
                Diberikan $P$ dan $Q=kP$, mencari $k$ sangat sukar. Inilah dasar
                ECDH dan ECEG.
            \end{alertblock}
        \end{column}
        \begin{column}{0.48\textwidth}
            \centering
            \bookfig[0.88\linewidth]{14}{sebaran-titik-gf11.png}
            \figsumber{Bab 14, \textit{Pengantar Kriptografi}.}
            \vspace{0.1cm}
            \begin{exampleblock}{Ukuran kunci}
                \footnotesize
                ECC 160 bit setara RSA 1024 bit; ECC 256 bit setara RSA
                3072 bit. Kunci pendek menghemat bandwidth dan memori, sehingga
                cocok untuk perangkat nirkabel.
            \end{exampleblock}
        \end{column}
    \end{columns}
\end{frame}

\begin{frame}{ECC: Satu Contoh Hitung}
    \scriptsize
    \begin{columns}[T,onlytextwidth]
        \begin{column}{0.50\textwidth}
            \begin{block}{Kurva $y^{2}\equiv x^{3}+x+6 \pmod{11}$}
                Kurva ini berorde 13. Titik $P(2,4)$ dan $Q(5,9)$ terletak
                pada kurva.
            \end{block}
            \begin{exampleblock}{Penjumlahan $P+Q$}
                $m \equiv (9-4)(5-2)^{-1} \equiv 5\cdot4 \equiv 9 \pmod{11}$.
                \par\vspace{1pt}
                $x_r \equiv 9^{2}-2-5=74\equiv8$;
                $y_r \equiv 9(2-8)-4=-58\equiv8$.
                \par\vspace{1pt}
                Jadi $P+Q=(8,8)$.
            \end{exampleblock}
            \begin{exampleblock}{Penggandaan $2P$}
                $m \equiv (3\cdot2^{2}+1)(2\cdot4)^{-1}\equiv13\cdot7\equiv3$;
                hasilnya $2P=(5,9)=Q$.
            \end{exampleblock}
        \end{column}
        \begin{column}{0.48\textwidth}
            \begin{block}{ECDH pada kurva kecil}
                Kurva $y^{2}\equiv x^{3}+2x+1\pmod5$ berorde 7 dengan basis
                $B(0,1)$. Alice memilih $a=2$, Bob $b=3$:
                \[
                    P_A=2B=(1,3), \qquad P_B=3B=(3,3).
                \]
                Kunci bersama:
                \[
                    K=2P_B=3P_A=(0,4).
                \]
            \end{block}
            \begin{alertblock}{Bacaan}
                \footnotesize
                Menghitung $kP$ mudah; mengembalikan $k$ dari $P$ dan $kP$
                sukar.
            \end{alertblock}
        \end{column}
    \end{columns}
\end{frame}

\begin{frame}{Knapsack Merkle--Hellman}
    \scriptsize
    \vspace{-0.1cm}
    \begin{columns}[T,onlytextwidth]
        \begin{column}{0.52\textwidth}
            \begin{block}{Gagasan}
                Knapsack mencari bit $b_i\in\{0,1\}$ dengan
                $\sum_i b_i w_i=M$. Versi umumnya \textit{NP-complete}.
                Barisan \textit{superincreasing} memenuhi
                $w_i>\sum_{j<i}w_j$, sehingga solusinya tunggal dan dapat
                dicari secara \textit{greedy} dalam $O(n)$.
            \end{block}
            \begin{block}{Menyembunyikan kemudahan}
                $w'_i=(w_i\cdot n)\bmod m$ dengan $m>\sum_i w_i$ dan
                $\operatorname{PBB}(n,m)=1$. Enkripsi
                $C=\sum_i b_i w'_i$; dekripsi
                $C'=(C\cdot n^{-1})\bmod m$, lalu \textit{greedy}.
            \end{block}
            \begin{exampleblock}{Contoh dari bab}
                $W_{priv}=\{2,3,6,13,27,52\}$, $m=105$, $n=31$ memberi
                $W_{pub}=\{62,93,81,88,102,37\}$ dan $n^{-1}=61$.
                Blok \code{110101} menjadi $280$; dekripsi $280$ memberi
                $C'=70$ lalu \code{110101}.
            \end{exampleblock}
        \end{column}
        \begin{column}{0.46\textwidth}
            \centering
            \bookfig[0.55\linewidth]{13}{ilustrasi-masalah-knapsack.png}
            \figsumber{Bab 13, \textit{Pengantar Kriptografi}.}
            \vspace{0.08cm}
            \begin{alertblock}{Mengapa patah}
                \footnotesize
                Shamir (1982) merekonstruksi barisan \textit{superincreasing}
                dari kunci publik dalam waktu polinomial. \kunci{Struktur
                rahasia} yang membangun sistem justru membocorkan jalan pintas.
            \end{alertblock}
        \end{column}
    \end{columns}
\end{frame}

\begin{frame}{Perbandingan Knapsack dan ECC}
    \scriptsize
    \begin{tabularx}{\linewidth}{@{}>{\raggedright\arraybackslash}p{2.8cm} >{\raggedright\arraybackslash}p{4.5cm} >{\raggedright\arraybackslash}X@{}}
        \toprule
        \textbf{Aspek} & \textbf{Knapsack Merkle--Hellman} & \textbf{ECC} \\
        \midrule
        Persoalan dasar & Himpunan bagian dengan jumlah tertentu & Logaritma diskret pada kurva eliptik (ECDLP) \\
        \addlinespace[1pt]
        Kunci publik & Barisan bobot $W_{pub}$ & Titik $Q=xB$ pada kurva \\
        \addlinespace[1pt]
        Operasi inti & Penjumlahan bobot dan modulo $m$ & Penjumlahan dan penggandaan titik \\
        \addlinespace[1pt]
        Ukuran cipherteks & Lebih besar daripada plainteks & Pasangan titik pada ECEG \\
        \addlinespace[1pt]
        Status & Patah sejak 1982 & Dipakai luas; kunci 256 bit setara RSA 3072 bit \\
        \addlinespace[1pt]
        Pelajaran & Persoalan sukar tidak menjamin sistem aman & Keamanan bergantung pada pemilihan kurva dan parameter \\
        \bottomrule
    \end{tabularx}
    \vspace{0.12cm}
    \begin{alertblock}{Satu pertanyaan pembeda}
        \footnotesize
        Knapsack menyembunyikan barisan mudah dengan transformasi rahasia;
        ECC tidak menyembunyikan struktur, tetapi memilih grup yang operasinya
        mudah dibalik satu arah dan sukar dibalik arah lain.
    \end{alertblock}
\end{frame}

\begin{frame}{Perbandingan Steganografi dan Kriptografi}
    \scriptsize
    \begin{columns}[T,onlytextwidth]
        \begin{column}{0.49\textwidth}
            \begin{block}{Gagasan}
                Steganografi menyembunyikan \kunci{keberadaan} pesan di dalam
                media sampul (\textit{cover}); hasilnya disebut
                \textit{stego-object}. Kriptografi menyembunyikan isi, tetapi
                keberadaan cipherteks tetap terlihat.
            \end{block}
            \begin{block}{LSB pada citra dan audio}
                Bit pesan menggantikan bit paling tidak berarti pada byte
                citra atau sampel audio. Perubahan satu tingkat hampir tidak
                terlihat atau terdengar. Semakin banyak bit dipakai, semakin
                besar kapasitas, tetapi semakin turun kualitas media.
            \end{block}
        \end{column}
        \begin{column}{0.49\textwidth}
            \begin{block}{Tabel pembanding}
                \begin{tabularx}{\linewidth}{@{}>{\raggedright\arraybackslash}p{2.3cm} >{\raggedright\arraybackslash}X >{\raggedright\arraybackslash}X@{}}
                    \toprule
                    \textbf{Aspek} & \textbf{Kriptografi} & \textbf{Steganografi} \\
                    \midrule
                    Yang disembunyikan & isi pesan & keberadaan pesan \\
                    \addlinespace[0pt]
                    Hasil terlihat & cipherteks & media tampak normal \\
                    \addlinespace[0pt]
                    Deteksi & ada, tetapi tidak terbaca & tidak diketahui ada \\
                    \bottomrule
                \end{tabularx}
            \end{block}
            \begin{block}{Steganalisis}
                \textit{Visual attack}, \textit{targeted steganalysis}, dan
                \textit{blind steganalysis}. Berhasil bila keputusannya lebih
                baik daripada menebak acak.
            \end{block}
        \end{column}
    \end{columns}
\end{frame}

% =====================================================================
\section{Latihan}
% =====================================================================

\begin{frame}{Latihan 1: Klasik dan One-Time Pad}
    \scriptsize
    \begin{columns}[T,onlytextwidth]
        \begin{column}{0.52\textwidth}
            \begin{enumerate}
                \setlength{\itemsep}{4pt}
                \item Dekripsikan cipherteks Caesar \code{KHOORZRUOG}.
                \item Enkripsikan \code{halo} dengan Affine $m=7$, $b=10$,
                      dan $n=26$.
                \item Pada cipherteks Vigen\`ere ditemukan kriptogram berulang
                      dengan jarak 15, 20, dan 30 huruf. Berapa panjang kunci
                      yang paling mungkin?
                \item Mengapa kunci one-time pad tidak boleh dipakai dua kali?
            \end{enumerate}
        \end{column}
        \begin{column}{0.46\textwidth}
            \begin{block}{Kunci jawaban}
                \begin{enumerate}
                    \setlength{\itemsep}{3pt}
                    \item \code{HELLO WORLD}; setiap huruf digeser tiga
                          ke belakang.
                    \item \code{HKJE}. Contoh: \code{h}$=7$ menjadi
                          $(7\cdot7+10)\bmod26=7=\code{h}$.
                    \item $\operatorname{PBB}(15,20,30)=5$, jadi panjang kunci
                          5 huruf.
                    \item Karena $C_1\oplus C_2=P_1\oplus P_2$; kunci saling
                          menghapus dan dua plainteks dapat dipisahkan.
                \end{enumerate}
            \end{block}
        \end{column}
    \end{columns}
\end{frame}

\begin{frame}{Latihan 2: Simetris Modern}
    \scriptsize
    \begin{columns}[T,onlytextwidth]
        \begin{column}{0.52\textwidth}
            \begin{enumerate}
                \setlength{\itemsep}{4pt}
                \item Hitung \code{10110010} $\oplus$ \code{11001010}.
                \item Berapa panjang kunci efektif DES, dan ke mana perginya
                      delapan bit sisanya?
                \item Terapkan \textit{MixColumns} AES pada kolom
                      $[\code{26}\ \code{7B}\ \code{BD}\ \code{43}]^{T}$.
                \item Pada mode CBC, satu bit blok cipherteks $C_3$ rusak.
                      Blok plainteks mana saja yang terpengaruh?
            \end{enumerate}
        \end{column}
        \begin{column}{0.46\textwidth}
            \begin{block}{Kunci jawaban}
                \begin{enumerate}
                    \setlength{\itemsep}{3pt}
                    \item \code{01111000}.
                    \item 56 bit; delapan bit pada posisi 8, 16, 24, 32, 40,
                          48, 56, dan 64 adalah paritas.
                    \item $[\code{3F}\ \code{4F}\ \code{F9}\ \code{2A}]^{T}$.
                    \item $P_3$ rusak seluruhnya dan $P_4$ berubah satu bit;
                          $P_5$ dan seterusnya pulih normal.
                \end{enumerate}
            \end{block}
        \end{column}
    \end{columns}
\end{frame}

\begin{frame}{Latihan 3: Kunci-Publik}
    \scriptsize
    \begin{columns}[T,onlytextwidth]
        \begin{column}{0.52\textwidth}
            \begin{enumerate}
                \setlength{\itemsep}{4pt}
                \item Dengan $p=47$, $q=71$, dan $e=17$, hitung
                      $\varphi(n)$ dan $d$.
                \item Dengan kunci publik $(79,3337)$, enkripsikan $m=16$.
                \item Pada Diffie--Hellman $p=353$, $g=3$, $a=97$, $b=233$,
                      hitung $A$, $B$, dan $K$.
                \item Pada ElGamal $p=2273$, $g=3$, $x=243$, hitung $y$.
                      Lalu enkripsikan $m=700$ dengan $k=1463$.
            \end{enumerate}
        \end{column}
        \begin{column}{0.46\textwidth}
            \begin{block}{Kunci jawaban}
                \begin{enumerate}
                    \setlength{\itemsep}{3pt}
                    \item $\varphi(n)=46\cdot70=3220$ dan
                          $d=17^{-1}\bmod3220=2273$.
                    \item $c=16^{79}\bmod3337=1493$; dekripsi mengembalikan 16.
                    \item $A=40$, $B=248$, dan $K=160$ dari kedua sisi.
                    \item $y=3^{243}\bmod2273=461$; cipherteksnya
                          $(a,b)=(1439,74)$.
                \end{enumerate}
            \end{block}
        \end{column}
    \end{columns}
\end{frame}

\begin{frame}{Latihan 4: Knapsack, ECC, dan Steganografi}
    \scriptsize
    \begin{columns}[T,onlytextwidth]
        \begin{column}{0.52\textwidth}
            \begin{enumerate}
                \setlength{\itemsep}{4pt}
                \item Apakah $W=\{3,7,14,29,60,118\}$ \textit{superincreasing}?
                      Pecahkan $M=96$ secara \textit{greedy}.
                \item Pada $y^{2}\equiv x^{3}+x+6\pmod{11}$, hitung
                      $(2,4)+(8,8)$.
                \item Kriptogram knapsack $C=280$ dengan $n^{-1}=61$ dan
                      $m=105$. Tentukan $C'$ dan blok bitnya.
                \item Sebutkan tiga jenis steganalisis dan tujuannya.
            \end{enumerate}
        \end{column}
        \begin{column}{0.46\textwidth}
            \begin{block}{Kunci jawaban}
                \begin{enumerate}
                    \setlength{\itemsep}{3pt}
                    \item Ya, karena setiap elemen melebihi jumlah semua
                          sebelumnya. Solusinya $7+29+60=96$, yaitu
                          \code{010110}.
                    \item Gradien $m=8$, sehingga hasilnya $(10,9)$.
                    \item $C'=(280\cdot61)\bmod105=70$; \textit{greedy} memberi
                          \code{110101}.
                    \item \textit{Visual}, \textit{targeted}, dan
                          \textit{blind}; tujuannya memutuskan apakah media
                          memuat pesan tersembunyi.
                \end{enumerate}
            \end{block}
        \end{column}
    \end{columns}
\end{frame}

% =====================================================================
\section{Penutup}
% =====================================================================

\begin{frame}{Enam Gagasan yang Harus Terbawa}
    \footnotesize
    \begin{columns}[T,onlytextwidth]
        \begin{column}{0.50\textwidth}
            \begin{block}{Yang harus diingat}
                \begin{enumerate}
                    \setlength{\itemsep}{1pt}
                    \item Keamanan bertumpu pada kunci, bukan pada kerahasiaan
                          algoritma.
                    \item Empat layanan menuntut mekanisme yang berbeda.
                    \item Cipher klasik runtuh karena statistik bahasa; cipher
                          modern dirancang agar statistik itu hilang.
                    \item Simetri menang di kecepatan; nir-simetri menang di
                          pengelolaan kunci.
                    \item Panjang kunci bukan penentu tunggal; struktur dan
                          persoalan matematis juga menentukan.
                    \item Steganografi menyembunyikan keberadaan; kriptografi
                          menyembunyikan isi.
                \end{enumerate}
            \end{block}
        \end{column}
        \begin{column}{0.48\textwidth}
            \begin{alertblock}{Kekeliruan yang paling sering muncul}
                \begin{itemize}
                    \setlength{\itemsep}{2pt}
                    \item Menganggap OTP praktis tanpa syarat kunci.
                    \item Memakai ECB pada data yang berulang.
                    \item Mengira RSA aman tanpa \textit{padding}.
                    \item Memakai Diffie--Hellman tanpa otentikasi.
                    \item Mengira knapsack aman karena persoalannya
                          \textit{NP-complete}.
                \end{itemize}
            \end{alertblock}
            \begin{exampleblock}{Cara belajar}
                \footnotesize
                Untuk setiap keluarga, hafalkan tiga hal: satuan pemrosesan,
                ukuran kunci, dan satu titik lemahnya.
            \end{exampleblock}
        \end{column}
    \end{columns}
\end{frame}

\begin{frame}{Penutup: Menuju Ujian Akhir Semester}
    \footnotesize
    \begin{columns}[T,onlytextwidth]
        \begin{column}{0.50\textwidth}
            \begin{block}{Cakupan UAS}
                Seluruh materi pertemuan 1--14: layanan dan serangan,
                matematika, cipher klasik dan OTP, cipher alir, cipher blok
                sampai AES, RSA, Diffie--Hellman, ElGamal, ECC, knapsack, serta
                steganografi dan steganalisis.
            \end{block}
            \begin{exampleblock}{Yang diperiksa}
                \footnotesize
                Definisi dan syarat setiap algoritma; ukuran kunci dan blok;
                satu contoh hitung; dan alasan mengapa sebuah algoritma
                ditinggalkan.
            \end{exampleblock}
        \end{column}
        \begin{column}{0.48\textwidth}
            \begin{alertblock}{Minggu 16 --- UAS}
                Pertemuan berikutnya adalah \kunci{UAS}, bukan pertemuan tatap
                muka. Tidak ada materi baru. Seluruh cakupan CPMK-1 dan CPMK-2
                yang telah dibahas pada empat belas pertemuan menjadi bahan
                ujian.
            \end{alertblock}
            \begin{block}{Sebelum masuk ruang ujian}
                \footnotesize
                Periksa kembali contoh angka RSA, Diffie--Hellman, ElGamal,
                knapsack, dan ECC. Kelimanya paling sering muncul sebagai soal
                hitung.
            \end{block}
        \end{column}
    \end{columns}
\end{frame}

\end{document}
